% !TeX root = part01.tex

\documentclass[../final.tex]{subfiles}

\begin{document}

% ============================================================
% Взаимодействие частиц с веществом
% ============================================================

\chapter{Взаимодействие частиц с веществом}

\rsubtitle{Ионизационные потери энергии}


% ============================================================
\section{Введение}
% ============================================================

% ------------------------------------------------------------
\subsection{Инвариантная масса}
% ------------------------------------------------------------

Для системы нескольких частиц удобно использовать
\rtxt{3}{инвариантную массу}, определяемую через суммарную энергию
и суммарный импульс:

\begin{equation*}
    M_{\mathrm{inv}}^2
    =
    \left(
        \sum_i E_i
    \right)^2
    -
    \left(
        \sum_i \mathbf p_i
    \right)^2.
\end{equation*}

Эта величина не зависит от выбора инерциальной системы отсчёта
и поэтому позволяет исследовать образование промежуточных состояний
и резонансов по измеренным энергиям и импульсам конечных частиц.

\begin{rbox}[left=18pt,right=18pt,top=14pt,bottom=10pt]{[]}
    \centering

    \includegraphics[
        width=0.5\linewidth
    ]{image/lecture_01/invariant_mass_peak.png}

    \captionsetup{
        font=footnotesize,
        labelfont=bf,
        skip=5pt
    }

    \captionof{figure}{
        Распределение по инвариантной массе системы
        $\pi^+\pi^-\pi^0$.
        Резонансное состояние проявляется в виде локального пика
        в распределении событий.
    }
    \label{fig:invariant_mass_peak}
\end{rbox}

Например, в реакции

\begin{equation*}
    p+p
    \rightarrow
    p+p+\pi^++\pi^-+\pi^0
\end{equation*}

можно построить распределение по инвариантной массе системы

\[
    \pi^+\pi^-\pi^0.
\]

Появление пика в таком распределении соответствует образованию
промежуточного резонансного состояния, которое затем распадается
на рассматриваемые частицы.


% ============================================================
\section{Ионизационные потери}
% ============================================================

При прохождении заряженной частицы через вещество происходит её
взаимодействие с атомами среды. Одним из основных процессов является
взаимодействие с атомными электронами, приводящее к возбуждению
и ионизации атомов.

В результате достаточно сильного столкновения электрон может быть
выбит из атома. Такой вторичный электрон называется
\rtxt{3}{$\delta$-электроном}.

\begin{rbox}[left=18pt,right=18pt,top=14pt,bottom=10pt]{[]}
    \centering

    \includegraphics[
        width=0.8\linewidth
    ]{image/lecture_01/delta_electron_scheme.png}

    \captionsetup{
        font=footnotesize,
        labelfont=bf,
        skip=5pt
    }

    \captionof{figure}{
        Ионизация вещества при прохождении заряженной частицы.
        В результате столкновения с атомным электроном может образоваться
        вторичный $\delta$-электрон с энергией $\varepsilon_\delta$.
    }
    \label{fig:delta_electron_scheme}
\end{rbox}


% ------------------------------------------------------------
\subsection{Максимальная энергия \texorpdfstring{$\delta$}{δ}-электрона}
% ------------------------------------------------------------

Пусть налетающая частица имеет массу $M$, энергию $E$ и импульс $p$.
Максимальная энергия, которую она может передать электрону массы $m_e$,
равна

\begin{equation*}
    \varepsilon_{\max}
    =
    2m_e
    \frac{p^2}{
        M^2
        +
        m_e^2
        +
        \dfrac{2Em_e}{c^2}
    }.
\end{equation*}

В зависимости от энергии налетающей частицы это выражение можно
упростить.

При нерелятивистском или слаборелятивистском движении,

\[
    \gamma \sim 1,
\]

получаем оценку

\begin{equation*}
    \varepsilon_{\max}[\mathrm{MeV}]
    \approx
    2(\gamma-1).
\end{equation*}

В промежуточной релятивистской области

\[
    1
    \ll
    \gamma
    \ll
    \frac{M}{2m_e}
\]

максимальная передаваемая энергия растёт приблизительно как

\begin{equation*}
    \varepsilon_{\max}[\mathrm{MeV}]
    \approx
    \gamma^2.
\end{equation*}

При очень больших энергиях,

\[
    \gamma
    \gg
    \frac{M}{2m_e},
\]

максимальная энергия $\delta$-электрона становится порядка энергии
самой налетающей частицы:

\[
    \varepsilon_{\max}
    \sim
    E.
\]


% ------------------------------------------------------------
\subsection{Вероятность взаимодействия частицы в веществе}
% ------------------------------------------------------------

Рассмотрим слой вещества толщины $\Delta x$. Если сечение
взаимодействия частицы с атомом равно $\sigma$, а концентрация атомов
равна $n$, то вероятность взаимодействия на малой толщине составляет

\begin{equation*}
    \Delta W
    =
    \sigma n\,\Delta x.
\end{equation*}

Концентрацию атомов вещества можно выразить через его плотность:

\begin{equation*}
    n
    =
    \frac{N_A}{A}\rho,
\end{equation*}

где $N_A$ --- число Авогадро, $A$ --- атомная масса,
$\rho$ --- плотность вещества.

Тогда

\begin{equation*}
    \Delta W
    =
    \sigma
    \Delta x
    \frac{N_A}{A}
    \rho.
\end{equation*}

Размерности соответствующих величин:

\begin{equation*}
    [\sigma]
    =
    \mathrm{cm^2},
    \qquad
    [n]
    =
    \mathrm{cm^{-3}}.
\end{equation*}

Вместо обычной геометрической толщины $\Delta x$ удобно использовать
\rtxt{3}{массовую толщину вещества}

\begin{equation*}
    \Delta\widetilde{x}
    =
    \rho\,\Delta x.
\end{equation*}

Она измеряется в

\begin{equation*}
    [\Delta\widetilde{x}]
    =
    \frac{\mathrm{g}}{\mathrm{cm^2}}.
\end{equation*}

Тогда вероятность взаимодействия можно записать как

\begin{equation*}
    \Delta W
    =
    \sigma
    \frac{N_A}{A}
    \Delta\widetilde{x}.
\end{equation*}

% ============================================================
\section{Рассеяние на атомных электронах}
% ============================================================

% ------------------------------------------------------------
\subsection{Формула Резерфорда}
% ------------------------------------------------------------

При кулоновском рассеянии заряженной частицы на электроне
дифференциальное сечение имеет резерфордовский вид

\begin{equation*}
    \frac{d\sigma}{d\Omega}
    =
    \frac{Z^2r_e^2}{4}
    \left(
        \frac{m_ec}{\beta_c p_c}
    \right)^2
    \frac{1}{
        \sin^4
        \left(
            \dfrac{\theta_c}{2}
        \right)
    }.
\end{equation*}

Здесь $r_e$ --- классический радиус электрона,
$p_c$ --- импульс в системе центра масс,

\[
    \beta_c
    =
    \frac{v_c}{c},
\]

а $\theta_c$ --- угол рассеяния в системе центра масс.

После перехода от угла рассеяния к энергии $\varepsilon$,
переданной электрону, сечение можно представить в виде

\begin{equation*}
    d\sigma
    =
    \frac{
        2\pi Z^2 r_e^2 m_ec^2
    }{
        \beta^2\varepsilon^2
    }
    d\varepsilon.
\end{equation*}

Здесь $\varepsilon$ --- энергия выбитого $\delta$-электрона.

Из этой зависимости видно, что вероятность передачи большой энергии
быстро уменьшается с ростом $\varepsilon$:

\[
    d\sigma
    \propto
    \frac{d\varepsilon}{\varepsilon^2}.
\]


% ------------------------------------------------------------
\subsection{Спектр \texorpdfstring{$\delta$}{δ}-электронов}
% ------------------------------------------------------------

Число $\delta$-электронов, образующихся при прохождении частицы
через вещество, определяется выражением

\begin{equation*}
    \frac{d^2n}{
        d\varepsilon\,
        d\widetilde{x}
    }
    =
    0.15354
    \frac{Z}{A}
    \frac{1}{
        \beta^2\varepsilon^2
    }.
\end{equation*}

Здесь $n$ --- число выбитых $\delta$-электронов,
$Z$ --- заряд ядра атома вещества,
$A$ --- атомная масса вещества,

\[
    \beta
    =
    \frac{v}{c},
\]

а $\varepsilon$ --- энергия $\delta$-электрона.

Основная зависимость спектра имеет вид

\begin{equation*}
    \frac{dn}{d\varepsilon}
    \propto
    \frac{1}{\varepsilon^2}.
\end{equation*}

Поэтому электронов с малой переданной энергией образуется значительно
больше, чем электронов с большой энергией.

\begin{rbox}[left=18pt,right=18pt,top=14pt,bottom=10pt]{[]}
    \centering

    \includegraphics[
        width=0.62\linewidth
    ]{image/lecture_01/delta_electron_spectrum.png}

    \captionsetup{
        font=footnotesize,
        labelfont=bf,
        skip=5pt
    }

    \captionof{figure}{
        Спектр $\delta$-электронов.
        Число выбитых электронов быстро уменьшается с ростом переданной энергии
        приблизительно по закону $1/\varepsilon^2$;
        максимальная возможная энергия ограничена величиной
        $\varepsilon_{\max}$.
    }
    \label{fig:delta_electron_spectrum}
\end{rbox}


% ------------------------------------------------------------
\subsection{Формула Мотта}
% ------------------------------------------------------------

Резерфордовская формула не учитывает релятивистскую спиновую структуру
электрона. Для учёта этого эффекта используется поправка Мотта:

\begin{equation*}
    \left(
        \frac{d\sigma}{d\Omega}
    \right)_{\!M}
    =
    \left(
        \frac{d\sigma}{d\Omega}
    \right)_{\!R}
    \left(
        1
        -
        \beta_c^2
        \sin^2
        \frac{\theta_c}{2}
    \right).
\end{equation*}

Здесь индекс $M$ соответствует формуле Мотта,
а индекс $R$ --- формуле Резерфорда.

После учёта этой поправки энергетический спектр $\delta$-электронов
принимает вид

\begin{equation*}
    \frac{d^2n}{
        d\varepsilon\,
        d\widetilde{x}
    }
    =
    0.15354
    \frac{Z}{A}
    \frac{1}{
        \beta^2\varepsilon^2
    }
    \left(
        1
        -
        \beta^2
        \frac{\varepsilon}{
            \varepsilon_{\max}
        }
    \right).
\end{equation*}

Здесь $\beta$ относится к налетающей частице.

При

\[
    \varepsilon
    \ll
    \varepsilon_{\max}
\]

поправочный множитель близок к единице, поэтому снова получается
приблизительная зависимость

\begin{equation*}
    \frac{d^2n}{
        d\varepsilon\,
        d\widetilde{x}
    }
    \propto
    \frac{1}{\varepsilon^2}.
\end{equation*}

При приближении энергии $\delta$-электрона к кинематическому пределу

\[
    \varepsilon
    \rightarrow
    \varepsilon_{\max}
\]

поправка становится существенной.


% ============================================================
\section{Удельные ионизационные потери}
% ============================================================

Пусть частица проходит через конечный слой вещества.

\begin{rbox}[left=18pt,right=18pt,top=14pt,bottom=10pt]{[]}
    \centering

    \includegraphics[
        width=0.72\linewidth
    ]{image/lecture_01/particle_through_matter.png}

    \captionsetup{
        font=footnotesize,
        labelfont=bf,
        skip=5pt
    }

    \captionof{figure}{
        Прохождение заряженной частицы через слой вещества толщины $\Delta x$.
        При движении через среду частица взаимодействует с атомами вещества
        и теряет часть своей энергии.
    }
    \label{fig:particle_through_matter}
\end{rbox}

Число столкновений с передачей энергии в интервале
от $\varepsilon$ до $\varepsilon+d\varepsilon$ определяется спектром

\[
    \frac{dn}{d\varepsilon}.
\]

Полное число взаимодействий получается интегрированием по допустимому
интервалу передаваемых энергий:

\begin{equation*}
    n
    =
    \int_{\varepsilon_{\min}}^{\varepsilon_{\max}}
    \frac{dn}{d\varepsilon}
    \,d\varepsilon.
\end{equation*}

При вычислении полной потерянной энергии каждое столкновение необходимо
дополнительно учитывать с весом, равным переданной энергии
$\varepsilon$.

Энергия, теряемая частицей при прохождении вещества, складывается
из энергий, переданных атомным электронам.

Для свободных электронов удельные энергетические потери имеют вид

\begin{equation*}
    \frac{dE}{dx}
    =
    K
    \frac{Z}{A}
    \frac{1}{\beta^2}
    \left[
        \ln
        \frac{
            \varepsilon_{\max}
        }{
            \varepsilon_{\min}
        }
        -
        \beta^2
        \frac{
            \varepsilon_{\max}
            -
            \varepsilon_{\min}
        }{
            \varepsilon_{\max}
        }
    \right].
\end{equation*}

Логарифмический член возникает при интегрировании энергетического
спектра. Поскольку при вычислении энергетических потерь число
столкновений умножается на переданную в одном столкновении энергию,
получаем

\begin{equation*}
    \int
    \varepsilon
    \frac{d\varepsilon}{\varepsilon^2}
    =
    \int
    \frac{d\varepsilon}{\varepsilon}.
\end{equation*}

Следовательно,

\begin{equation*}
    \int_{\varepsilon_{\min}}^{\varepsilon_{\max}}
    \frac{d\varepsilon}{\varepsilon}
    =
    \ln
    \frac{
        \varepsilon_{\max}
    }{
        \varepsilon_{\min}
    }.
\end{equation*}

Второй член в выражении для потерь связан с поправкой Мотта

\[
    1
    -
    \beta^2
    \frac{\varepsilon}{
        \varepsilon_{\max}
    }.
\]

В результате зависимость ионизационных потерь от скорости частицы
имеет характерную форму.

\begin{rbox}[left=18pt,right=18pt,top=14pt,bottom=10pt]{[]}
    \centering

    \includegraphics[
        width=0.72\linewidth
    ]{image/lecture_01/ionization_losses_curve.png}

    \captionsetup{
        font=footnotesize,
        labelfont=bf,
        skip=5pt
    }

    \captionof{figure}{
        Качественная зависимость удельных ионизационных потерь от скорости частицы.
        При малых скоростях потери велики, затем уменьшаются до минимума,
        после чего начинается медленный логарифмический рост.
    }
    \label{fig:ionization_losses_curve}
\end{rbox}

При малых и умеренных скоростях существенную роль играет множитель

\[
    \frac{1}{\beta^2},
\]

поэтому с увеличением скорости частицы потери уменьшаются.

При больших релятивистских энергиях начинает проявляться
логарифмический рост.


% ------------------------------------------------------------
\subsection{Формула Бете--Блоха}
% ------------------------------------------------------------

До сих пор электроны рассматривались как свободные. В реальном
веществе электроны связаны в атомах, поэтому необходимо учитывать
характерную энергию их возбуждения и ионизации.

Удельные ионизационные потери тяжёлой заряженной частицы описываются
\rtxt{3}{формулой Бете--Блоха}:

\begin{equation*}
    -
    \left(
        \frac{dE}{dx}
    \right)
    =
    Kz^2
    \frac{Z}{A}
    \frac{1}{\beta^2}
    \left[
        \frac{1}{2}
        \ln
        \left(
            \frac{
                2m_ec^2
                \beta^2
                \gamma^2
                \varepsilon_{\max}
            }{
                I^2
            }
        \right)
        -
        \beta^2
        -
        \frac{\delta}{2}
    \right].
\end{equation*}

Здесь $z$ --- заряд налетающей частицы,
$Z$ --- заряд ядра атома вещества,
$A$ --- атомная масса вещества,
$\beta=v/c$ --- скорость налетающей частицы,
$\gamma$ --- её лоренц-фактор,
$m_e$ --- масса электрона,
$\varepsilon_{\max}$ --- максимальная энергия, которую частица может
передать одному электрону,
$I$ --- средний ионизационный потенциал вещества,
$\delta$ --- поправка на эффект плотности.

Константа $K$ приблизительно равна

\begin{equation*}
    K
    \simeq
    0.307\,
    \frac{\mathrm{MeV}}{\mathrm{g/cm^2}},
\end{equation*}

или, что эквивалентно,

\begin{equation*}
    K
    \simeq
    0.307\,
    \mathrm{MeV\,cm^2/g}.
\end{equation*}

Средний ионизационный потенциал $I$ учитывает тот факт, что электроны
вещества связаны в атомных оболочках и для их возбуждения или
ионизации необходимо передать конечную энергию.

Поправка $\delta$ учитывает эффект плотности среды.


% ============================================================
\section{Флуктуации ионизационных потерь}
% ============================================================

Энергетические потери частицы в веществе не являются строго
фиксированными. Даже если частицы имеют одинаковую начальную энергию
и проходят через одинаковую толщину вещества, в каждом отдельном
событии потерянная энергия может различаться.

Причиной является статистический характер отдельных столкновений
с атомными электронами.

Для характеристики ширины распределения энергетических потерь
используется величина

\begin{equation*}
    \sigma_{\varepsilon}^2
    =
    \int
    n_i\varepsilon^2
    \,d\varepsilon.
\end{equation*}


% ------------------------------------------------------------
\subsection{Распределение Ландау}
% ------------------------------------------------------------

Для достаточно тонких слоёв вещества распределение энергетических
потерь является существенно несимметричным и описывается
\rtxt{3}{распределением Ландау}.

\begin{rbox}[left=18pt,right=18pt,top=14pt,bottom=10pt]{[]}
    \centering

    \includegraphics[
        width=0.72\linewidth
    ]{image/lecture_01/landau_distribution.png}

    \captionsetup{
        font=footnotesize,
        labelfont=bf,
        skip=5pt
    }

    \captionof{figure}{
        Распределение Ландау для флуктуаций ионизационных потерь.
        Распределение имеет выраженный максимум и длинный правый хвост,
        связанный с редкими столкновениями с большой передачей энергии;
        при больших $\varepsilon$ хвост убывает приблизительно как
        $1/\varepsilon^2$.
    }
    \label{fig:landau_distribution}
\end{rbox}

Длинный правый хвост распределения связан с редкими столкновениями,
в которых налетающая частица передаёт одному электрону сравнительно
большую энергию.

Поэтому для такого распределения необходимо различать среднюю
и наиболее вероятную потерю энергии.


% ------------------------------------------------------------
\subsection{Наиболее вероятная потеря энергии}
% ------------------------------------------------------------

Наиболее вероятное значение энергетических потерь обозначим как \( \Delta E_{\mathrm{\text{н.в.}}}. \) 

Для него записывается выражение

\begin{equation*}
    \Delta E_{\mathrm{\text{н.в.}}}
    =
    \xi
    \left[
        \ln
        \left(
            \frac{
                2m_ec^2
                \beta^2
                \gamma^2
                \xi
            }{
                I^2
            }
        \right)
        -
        \beta^2
        -
        \delta
        +
        j
    \right].
\end{equation*}

Постоянная $j$ приблизительно равна

\[
    j
    \simeq
    0.2.
\]

Параметр $\xi$ определяется как

\begin{equation*}
    \xi
    =
    0.154
    \frac{Z}{A}
    \frac{X}{\beta^2}.
\end{equation*}

Здесь $X$ характеризует толщину пройденного вещества в массовых
единицах.

Поправка

\[
    \delta
\]

связана с эффектом плотности среды.

\end{document}